\documentclass[11pt]{article}

% Packages
\usepackage[utf8]{inputenc}
\usepackage[margin=1in]{geometry}
\usepackage{amsmath,amssymb}
\usepackage{booktabs}
\usepackage{graphicx}
\usepackage{float}
\usepackage{hyperref}
\usepackage{caption}
\usepackage{subcaption}
\usepackage{natbib}
\usepackage{setspace}
\onehalfspacing

% Title
\title{Descriptive Data Analysis\\
\large 1907 Panic Paper}
\author{Generated by Warp}
\date{\today}

\begin{document}

\maketitle

\tableofcontents
\newpage

% =============================================================================
\section{Introduction}
% =============================================================================

This document presents descriptive statistics and exploratory data analysis for the 1907 Panic Paper project. The analysis covers the following datasets:

\begin{itemize}
    \item Census data (county-level demographics)
    \item Electoral data (presidential election results 1904--1912)
    \item Bank failure data (digitized records)
    \item Newspaper data (Chronicling America)
\end{itemize}

% =============================================================================
\section{Data Overview}
% =============================================================================

\subsection{Electoral Data}

The electoral dataset contains county-level presidential election returns for the elections of 1904, 1908, and 1912. Key variables include:

\begin{itemize}
    \item Republican vote share (percentage)
    \item Democratic vote share (percentage)
    \item Socialist vote share (percentage)
    \item Swing variables measuring change between elections
\end{itemize}

% Include summary statistics table (generated by R)
% \input{../Data/data_outputs/descriptive_reports/elections_summary.tex}

\subsection{Bank Failure Data}

The bank failure dataset contains digitized records of bank failures during and after the 1907 panic.

% \input{../Data/data_outputs/descriptive_reports/bank_failures_summary.tex}

\subsection{Census Data}

County-level demographic data from the U.S. Census, including population, literacy rates, and economic indicators.

% \input{../Data/data_outputs/descriptive_reports/census_summary.tex}

% =============================================================================
\section{Summary Statistics}
% =============================================================================

\subsection{Electoral Variables}

Table~\ref{tab:elections_summary} presents summary statistics for the key electoral variables.

\begin{table}[H]
\centering
\caption{Summary Statistics: Electoral Data}
\label{tab:elections_summary}
\begin{tabular}{lrrrrr}
\toprule
Variable & N & Mean & SD & Min & Max \\
\midrule
\multicolumn{6}{l}{\textit{1904 Election}} \\
Republican \% & -- & -- & -- & -- & -- \\
Democratic \% & -- & -- & -- & -- & -- \\
\midrule
\multicolumn{6}{l}{\textit{1908 Election}} \\
Republican \% & -- & -- & -- & -- & -- \\
Democratic \% & -- & -- & -- & -- & -- \\
\midrule
\multicolumn{6}{l}{\textit{Swing Variables (1904--1908)}} \\
Rep. Swing & -- & -- & -- & -- & -- \\
Anti-Rep. Swing & -- & -- & -- & -- & -- \\
Socialist Swing & -- & -- & -- & -- & -- \\
\bottomrule
\end{tabular}
\vspace{0.5em}
\footnotesize\textit{Note:} Fill in values from R output or use \texttt{stargazer} to generate.
\end{table}

% =============================================================================
\section{Missing Data Analysis}
% =============================================================================

Understanding patterns of missing data is crucial for valid inference.

\subsection{Missing Data by Variable}

% Include missing data summary

\subsection{Complete Cases}

% Report number of complete cases

% =============================================================================
\section{Distributions}
% =============================================================================

\subsection{Electoral Swing Variables}

Figure~\ref{fig:distributions} shows the distributions of key swing variables.

\begin{figure}[H]
\centering
% \includegraphics[width=0.9\textwidth]{../Data/data_outputs/descriptive_reports/elections_distributions.pdf}
\caption{Distributions of Electoral Swing Variables (1904--1908)}
\label{fig:distributions}
\end{figure}

\subsection{Outlier Detection}

Using the interquartile range (IQR) method with a multiplier of 1.5, we identify potential outliers in the data.

% Include outlier table

% =============================================================================
\section{Correlations}
% =============================================================================

\subsection{Electoral Variables}

Figure~\ref{fig:correlations} presents the correlation matrix for electoral variables.

\begin{figure}[H]
\centering
% \includegraphics[width=0.7\textwidth]{../Data/data_outputs/descriptive_reports/elections_correlations.pdf}
\caption{Correlation Matrix: Electoral Variables}
\label{fig:correlations}
\end{figure}

% =============================================================================
\section{Key Findings}
% =============================================================================

\begin{enumerate}
    \item \textbf{Republican Vote Decline:} [Describe patterns in anti-Republican swing]
    
    \item \textbf{Socialist Vote Increase:} [Describe patterns in socialist swing]
    
    \item \textbf{Geographic Patterns:} [Describe regional variation]
    
    \item \textbf{Data Quality:} [Note any issues with missing data or outliers]
\end{enumerate}

% =============================================================================
\section{Next Steps}
% =============================================================================

Based on this exploratory analysis, the following next steps are recommended:

\begin{itemize}
    \item Merge electoral data with bank failure data at the county level
    \item Control for pre-existing trends using 1900 election data
    \item Consider spatial autocorrelation in the analysis
    \item Address missing data through multiple imputation if necessary
\end{itemize}

% =============================================================================
% References (if needed)
% =============================================================================

% \bibliographystyle{apalike}
% \bibliography{references}

\end{document}
